% =============================================================
%  Tugas Individu 06 — Strategi Fusi Multimodal
%  IF25-40304 · Pembelajaran Mesin Multimodal · ITERA
% =============================================================
\documentclass[12pt, a4paper]{article}

\usepackage{tugas-style}

\renewcommand{\assignmentnumber}{06}
\renewcommand{\doctitle}{Tugas Individu 06}
\renewcommand{\duedate}{\color{ctpBlue}\faCalendarCheck\ \ Tenggat: sesuai jadwal kelas \faCalendarCheck}

\begin{document}

% ── HEADER ──────────────────────────────────────────────────
\makeassignmentheader

\hrule width0.3\textwidth
\revisioninfo{%
  \vhEntry{0.1}{16 September 2026}{I.W.W.}{Draf inisial tugas.}
}
\hrule
\vskip .3in

\tableofcontents
\clearpage

% =============================================================
\section{Tujuan Pembelajaran}\label{sec:objectives}

Setelah menyelesaikan tugas ini, mahasiswa mampu:
\begin{objectives}
  \item Membuktikan secara empiris bahwa satu modalitas saja tidak cukup, melalui pengukuran \emph{baseline} unimodal.
  \item Mengimplementasikan \emph{Early Fusion} pada tingkat fitur hasil encoder.
  \item Mengimplementasikan \emph{Late Fusion} pada tingkat keputusan.
  \item Mengimplementasikan \emph{Intermediate Fusion} pada lapisan tengah.
  \item Membandingkan kinerja dan ketangguhan ketiga strategi, lalu memilih serta mempertanggungjawabkan pilihannya.
\end{objectives}

% =============================================================
\section{Catatan dan Batasan}\label{sec:notes}
\vspace{-.1cm}

\begin{minipage}{\linewidth}
\begin{bclogo}[couleur=ctpBase, arrondi=0.1, logo=\bccrayon, ombre=false, couleurBord=ctpSurface, epBord=0.6]{Perhatikan hal berikut:}
  \begin{coloredPen}
    \item Tugas ini bersifat \textbf{individu}; tidak diperkenankan bekerja sama atau menyalin jawaban.
    \item \textbf{Encoder dibekukan.} Seluruh encoder modalitas sudah disediakan dan \textbf{tidak boleh dilatih}. Fokus tugas adalah lapisan fusi.
    \item Data bersifat \textbf{sintetis} dan dibangkitkan secara programatik dengan \emph{seed} tetap, sehingga hasil dapat direproduksi tanpa mengunduh dataset.
    \item Laporkan selalu \textbf{baseline unimodal} pada metrik yang sama. Angka akurasi fusi tidak bermakna tanpa pembanding.
    \item Gunakan \code{torch} untuk model dan \code{scikit-learn} untuk pembagian data serta metrik.
  \end{coloredPen}
\end{bclogo}
\end{minipage}

\begin{minipage}{\linewidth}
\begin{bclogo}[couleur=ctpMantle, arrondi=0.1, logo=\bcinfo, ombre=false, couleurBord=ctpSurface, epBord=0.6]{Ekspektasi Hasil}
  Dengan encoder \code{tanh} yang beku, hasil yang wajar adalah: kedua model unimodal
  sekitar \textbf{50\%}, \emph{Early Fusion} di atas \textbf{95\%},
  \emph{Intermediate Fusion} sedikit di bawahnya, dan \emph{Late Fusion}
  sekitar \textbf{85--90\%}.

  Urutan ini berlaku untuk label yang bersifat \textbf{aditif} (seperti pada tugas
  ini). Untuk label yang menuntut \textbf{interaksi} antar modalitas (misalnya
  paritas/XOR), \emph{Late Fusion} akan jatuh ke sekitar 50\% sementara
  \emph{Early} dan \emph{Intermediate} tetap tinggi.

  Bila hasil Anda menyimpang jauh (misalnya Early di bawah 80\%), periksa: encoder
  memakai \code{ReLU} sehingga dimensi informatif mati (gunakan \code{Tanh});
  label salah digabung; atau encoder tanpa sengaja ikut terlatih.
\end{bclogo}
\end{minipage}

% =============================================================
\section{Perangkat Lunak dan Peralatan yang Diperlukan}\label{sec:requiredsw}

\begin{itemize}[itemsep=2pt,parsep=0pt,topsep=2pt,partopsep=2pt]
  \item[\color{ctpBlue}\faPython] \textbf{Python} 3.8 atau lebih baru.
  \item[\color{ctpBlue}\faFileCode] \textbf{Paket wajib:} \code{numpy}, \code{torch}, \code{scikit-learn}, \code{matplotlib}.
  \item[\color{ctpBlue}\faDownload] \textbf{Data:} dibangkitkan sendiri (sintetis) --- tidak ada unduhan dataset.
  \item[\color{ctpBlue}\faFileCode] \textbf{Perangkat:} CPU sudah memadai; GPU tidak diperlukan.
\end{itemize}

\begin{codewindow}{setup.sh}
\begin{lstlisting}[style=pycode]
pip install numpy torch scikit-learn matplotlib
\end{lstlisting}
\end{codewindow}

\begin{minipage}{\linewidth}
\begin{bclogo}[couleur=ctpMantle, arrondi=0.1, logo=\bcinfo, ombre=false, couleurBord=ctpSurface, epBord=0.6]{Batasan Cakupan}
  Tugas ini memakai data multimodal \textbf{sintetis} yang dirancang agar setiap modalitas
  hanya membawa \textbf{separuh} informasi. Akibatnya, satu modalitas saja tidak pernah
  melampaui akurasi kebetulan, sedangkan penggabungan yang tepat dapat mencapainya.
  Konstruksi ini membuat perbedaan Early, Late, dan Intermediate Fusion terlihat jelas,
  tanpa memerlukan dataset eksternal.
\end{bclogo}
\end{minipage}

% =============================================================
\section{Pernyataan Masalah}\label{sec:intro}

Sebuah sistem diminta memprediksi \textbf{empat kelas emosi} dari dua modalitas: fitur
visual (wajah) dan fitur audio (suara). Setiap modalitas hanya membawa satu dari dua
faktor laten yang menentukan kelas, sehingga:

\begin{enumerate}
  \item Model unimodal visual paling tinggi mencapai akurasi sekitar 50\% --- ia tahu separuh jawaban.
  \item Model unimodal audio juga paling tinggi sekitar 50\%.
  \item Hanya penggabungan yang tepat dapat melampaui 95\%.
\end{enumerate}

Anda akan menguji ketiga strategi fusi pada kondisi ini, kemudian memilih yang paling
tepat beserta alasannya. Rincian pengerjaan diberikan pada bagian berikutnya.

% =============================================================
\section{Persyaratan}\label{sec:requirements}

Perhatikan setiap persyaratan pada soal-soal berikut.

% ------------------------------------------------------------
\subsection{Soal 1 --- Data Sintetis dan Baseline Unimodal}
\label{sec:soal1}

Sebelum menggabungkan apa pun, kita harus tahu seberapa jauh satu modalitas saja mampu
menjangkau. Inilah \emph{baseline} yang akan dibandingkan sepanjang tugas.

\subsubsection{Instruksi}
\label{sec:soal1-instruksi}

\textbf{(a)} Bangkitkan data sintetis dengan dua faktor laten biner (valensi dari wajah,
gairah dari suara). Label adalah gabungan keduanya sehingga menghasilkan empat kelas.

\textbf{(b)} Tampilkan distribusi kelas pada data. Pastikan keempat kelas berimbang.

\textbf{(c)} Latih dua model unimodal terpisah: satu hanya menerima fitur visual, satu
hanya menerima fitur audio. Laporkan akurasi masing-masing di data uji.

\textbf{(d)} Jawab dalam 3--5 kalimat: mengapa akurasi kedua model unimodal berhenti di
sekitar 50\% dan tidak dapat melampauinya?
\subsubsection{Kerangka Kode}
\label{sec:soal1-kode}

\begin{codewindow}{soal1\_data.py}
\begin{lstlisting}[style=pycode]
import numpy as np
import torch
import torch.nn as nn

SEED, N, DIM = 42, 4000, 16
rng = np.random.default_rng(SEED)

def buat_data(n=N):
    """Dua faktor laten: visual (valensi) dan audio (gairah)."""
    z_visual = rng.integers(0, 2, size=n)   # separuh jawaban
    z_audio  = rng.integers(0, 2, size=n)   # separuh jawaban

    visual = rng.normal(0, 1.0, size=(n, DIM)).astype(np.float32)
    audio  = rng.normal(0, 1.0, size=(n, DIM)).astype(np.float32)

    # Dimensi 0 menandai faktor laten; sisanya derau tak informatif
    visual[:, 0] = z_visual * 2.0 - 1.0
    audio[:, 0]  = z_audio  * 2.0 - 1.0

    label = z_visual * 2 + z_audio          # 4 kelas gabungan
    return visual, audio, label

visual, audio, label = buat_data()
print("visual:", visual.shape, "| audio:", audio.shape)
print("distribusi kelas:", np.bincount(label))

def buat_encoder(fan_in=DIM, fan_out=32, seed=0):
    """Encoder BEKU dengan aktivasi tanh.

    Catatan: gunakan tanh, bukan ReLU. ReLU dapat mematikan
    dimensi yang membawa faktor laten, sehingga informasi hilang
    sebelum sempat digabungkan.
    """
    g = torch.Generator().manual_seed(seed)
    enc = nn.Sequential(
        nn.Linear(fan_in, fan_out), nn.Tanh(),
        nn.Linear(fan_out, fan_out), nn.Tanh(),
    )
    with torch.no_grad():                   # inisialisasi terkendali
        for p in enc.parameters():
            if p.dim() > 1:
                p.copy_(torch.randn(p.shape, generator=g) / p.shape[1] ** 0.5)
            else:
                p.zero_()
    enc.eval()
    for p in enc.parameters():
        p.requires_grad_(False)             # bekukan
    return enc

ENC_VISUAL = buat_encoder(seed=1)
ENC_AUDIO  = buat_encoder(seed=2)
print("parameter encoder visual:", sum(p.numel() for p in ENC_VISUAL.parameters()))
\end{lstlisting}
\end{codewindow}

% ------------------------------------------------------------
\subsection{Soal 2 --- Early Fusion}
\label{sec:soal2}

Early Fusion menggabungkan fitur tepat setelah encoder, lalu meneruskannya ke satu model
bersama. Semua pembelajaran terjadi pada model gabungan tersebut.

\subsubsection{Instruksi}
\label{sec:soal2-instruksi}

\textbf{(a)} Bentuk vektor gabungan $\mathbf{z} = [\,f_A(\mathbf{x}_A) ; f_B(\mathbf{x}_B)\,]$
dari keluaran kedua encoder beku.

\textbf{(b)} Latih satu model bersama (MLP) di atas $\mathbf{z}$, lalu laporkan akurasinya.

\textbf{(c)} Bandingkan dengan baseline unimodal Soal 1. Apakah fusi menambah nilai?

\begin{codewindow}{soal2\_early.py}
\begin{lstlisting}[style=pycode]
class EarlyFusion(nn.Module):
    """Gabungkan fitur hasil encoder, lalu satu model bersama."""
    def __init__(self, dim_enc=32, n_kelas=4):
        super().__init__()
        self.head = nn.Sequential(
            nn.Linear(dim_enc * 2, 64), nn.Tanh(),
            nn.Linear(64, n_kelas),
        )

    def forward(self, x_visual, x_audio):
        with torch.no_grad():               # encoder beku
            f_a = ENC_VISUAL(x_visual)
            f_b = ENC_AUDIO(x_audio)
        z = torch.cat([f_a, f_b], dim=1)    # konkatenasi di lapisan awal
        return self.head(z)

model = EarlyFusion()
akurasi = latih_dan_uji(model, visual, audio, label)
print(f"Early Fusion accuracy: {akurasi:.4f}")
\end{lstlisting}
\end{codewindow}

% ------------------------------------------------------------
\subsection{Soal 3 --- Late Fusion}
\label{sec:soal3}

Late Fusion melatih model terpisah per modalitas, lalu menggabungkan \emph{keputusan}
masing-masing, bukan fiturnya.

\subsubsection{Instruksi}
\label{sec:soal3-instruksi}

\textbf{(a)} Latih dua model unimodal terpisah di atas fitur encoder, masing-masing
menghasilkan skor keyakinan untuk keempat kelas.

\textbf{(b)} Gabungkan kedua skor dengan rata-rata terbobot, lalu laporkan akurasinya.

\textbf{(c)} Jawab dalam 3--5 kalimat: mengapa strategi ini \textbf{masih dapat
melampaui baseline unimodal} ($\sim$50\%) meskipun encoder tidak boleh dilatih?
Jelaskan peran rata-rata terbobot, lalu sebutkan satu kondisi di mana Late Fusion
\textbf{akan gagal}.

\begin{codewindow}{soal3\_late.py}
\begin{lstlisting}[style=pycode]
class LateFusion(nn.Module):
    """Dua kepala unimodal; keputusan digabung di akhir."""
    def __init__(self, dim_enc=32, n_kelas=4):
        super().__init__()
        self.head_a = nn.Sequential(nn.Linear(dim_enc, n_kelas))
        self.head_b = nn.Sequential(nn.Linear(dim_enc, n_kelas))
        self.w_a = nn.Parameter(torch.tensor(0.5))   # bobot modalitas
        self.w_b = nn.Parameter(torch.tensor(0.5))

    def forward(self, x_visual, x_audio):
        with torch.no_grad():
            f_a = ENC_VISUAL(x_visual)
            f_b = ENC_AUDIO(x_audio)
        p_a = torch.softmax(self.head_a(f_a), dim=1)
        p_b = torch.softmax(self.head_b(f_b), dim=1)
        return self.w_a * p_a + self.w_b * p_b      # fusi keputusan
\end{lstlisting}
\end{codewindow}

% ------------------------------------------------------------
\clearpage
\subsection{Soal 4 --- Intermediate Fusion}
\label{sec:soal4}

Intermediate Fusion menggabungkan representasi pada lapisan tengah: setelah tiap modalitas
ditransformasi sedikit, tetapi sebelum keputusan akhir diambil.

\subsubsection{Instruksi}
\label{sec:soal4-instruksi}

\textbf{(a)} Tambahkan transformasi kecil per modalitas (lapisan tengah), lalu gabungkan
hasilnya. Anda boleh memakai konkatenasi atau \emph{gating}.

\textbf{(b)} Latih dan laporkan akurasinya.

\textbf{(c)} Bandingkan dengan Early dan Late Fusion. Anda akan melihat Intermediate
sedikit \emph{di bawah} Early karena encoder bersifat beku: lapisan tengah tambahan hanya
menambah transformasi sebelum informasi dari kedua modalitas bertemu. Tuliskan analisis
Anda dalam 4--6 kalimat, dan sebutkan bagaimana hasilnya akan berbeda bila encoder boleh
ikut dilatih (\emph{fine-tuning}).

\begin{codewindow}{soal4\_intermediate.py}
\begin{lstlisting}[style=pycode]
class IntermediateFusion(nn.Module):
    """Fusi pada lapisan tengah, setelah transformasi per modalitas."""
    def __init__(self, dim_enc=32, dim_mid=32, n_kelas=4):
        super().__init__()
        self.mid_a = nn.Sequential(nn.Linear(dim_enc, dim_mid), nn.Tanh())
        self.mid_b = nn.Sequential(nn.Linear(dim_enc, dim_mid), nn.Tanh())
        self.gate = nn.Sequential(nn.Linear(dim_mid * 2, dim_mid), nn.Sigmoid())
        self.head = nn.Sequential(
            nn.Linear(dim_mid, 64), nn.ReLU(),
            nn.Linear(64, n_kelas),
        )

    def forward(self, x_visual, x_audio):
        with torch.no_grad():
            f_a = ENC_VISUAL(x_visual)
            f_b = ENC_AUDIO(x_audio)
        h_a, h_b = self.mid_a(f_a), self.mid_b(f_b)
        g = self.gate(torch.cat([h_a, h_b], dim=1))
        h = g * h_a + (1 - g) * h_b          # gerbang adaptif
        return self.head(h)
\end{lstlisting}
\end{codewindow}

% ------------------------------------------------------------
\subsection{Soal 5 --- Analisis Kritis dan Uji Modalitas Hilang}
\label{sec:soal5}

Jawab pertanyaan berikut dalam bentuk paragraf singkat (masing-masing 5--8 kalimat).

\textbf{(a) Perbandingan dan Pemilihan}

Sajikan akurasi ketiga strategi beserta baseline unimodal dalam satu tabel. Pilih satu
strategi untuk sistem ini dan pertanggungjawabkan pilihan Anda berdasarkan akurasi,
kompleksitas, dan kemudahan implementasi.

\textbf{(b) Uji Modalitas Hilang}

Ulangi pengujian dengan salah satu modalitas dinolkan (misalkan fitur audio diisi nol).
Catat penurunan akurasi tiap strategi. Strategi mana yang paling tangguh, dan mengapa?

\begin{codewindow}{soal5\_hilang.py}
\begin{lstlisting}[style=pycode]
# Nolkan satu modalitas -> ukur ketangguhan tiap strategi
audio_kosong = torch.zeros_like(torch.as_tensor(audio))

for nama, model in [("Early", m_early),
                    ("Late", m_late),
                    ("Intermediate", m_inter)]:
    akurasi_normal = uji(model, visual, audio, label)
    akurasi_hilang = uji(model, visual, audio_kosong, label)
    print(f"{nama:14s} normal={akurasi_normal:.3f} "
          f"audio-hilang={akurasi_hilang:.3f}")
\end{lstlisting}
\end{codewindow}

\textbf{(c) Kaitan dengan Teori}

Hubungkan hasil Anda dengan dua gagasan dari perkuliahan: (i) masalah \emph{redundansi},
\emph{konflik}, dan \emph{dominasi modalitas}; serta (ii) peringatan bahwa penambahan
modalitas \textbf{tidak otomatis} meningkatkan kinerja. Apakah data sintetis pada tugas
ini termasuk kasus yang menguntungkan bagi fusi? Jelaskan.

% =============================================================
\clearpage
\section{Kriteria Penilaian}\label{sec:evaluation}

Tugas Anda akan dinilai berdasarkan kriteria berikut:

\renewcommand{\arraystretch}{1.5}
\begin{center}\small
\begin{tabular}{|p{11.5cm}|c|}
  \hline
  \thead{\color{ctpBlue} Kriteria} & \thead{\color{ctpBlue}Bobot} \\
  \hline
  Kebenaran teknis --- encoder tetap beku, kode berjalan tanpa error, ketiga strategi diimplementasikan dengan tepat. & 35\% \\
  \hline
  Kedalaman analisis --- penjelasan menghubungkan hasil dengan konsep fusi, baseline unimodal, dan masalah klasik fusi. & 30\% \\
  \hline
  Eksperimen --- pengukuran baseline, perbandingan ketiga strategi, dan uji modalitas hilang. & 20\% \\
  \hline
  Kualitas presentasi --- tabel hasil rapi, laporan terstruktur, dan interpretasi jelas. & 15\% \\
  \hline
  \textbf{Total} & \textbf{100\%} \\
  \hline
\end{tabular}
\end{center}

% =============================================================
\section{Apa yang Dikumpulkan}\label{sec:submit}

Seluruh artefak tugas dikemas ke dalam \textbf{satu berkas arsip} dengan format penamaan:
\begin{center}
  \code{mgg06\_tugas\_nim.zip}
\end{center}

dengan \code{nim} diganti oleh Nomor Induk Mahasiswa (NIM) masing-masing tanpa spasi.

\textbf{Artefak yang dapat disertakan di dalam arsip} (dapat lebih dari satu jenis berkas, disesuaikan dengan bentuk penyelesaian Anda):
\begin{borderedsquare}
  \item \code{mgg06\_tugas\_NIM.ipynb} --- apabila penyelesaian dikerjakan menggunakan Jupyter Notebook (termasuk Google Colab). Notebook \textbf{wajib telah dijalankan} sehingga seluruh \emph{output} sel tersimpan di dalam berkas.
  \item \code{mgg06\_tugas\_NIM.py} --- apabila penyelesaian dikerjakan dalam bentuk skrip Python.
  \item \code{mgg06\_tugas\_NIM.pdf} --- apabila analisis ditulis dalam dokumen teks terpisah.
  \item \code{mgg06\_tugas\_xx.jpg} --- apabila terdapat plot atau gambar pendukung tambahan; \code{xx} merupakan nomor urut visualisasi (misalnya \code{01}, \code{02}, dan seterusnya).
\end{borderedsquare}

Setiap berkas menggunakan awalan \code{mgg06\_tugas\_NIM} yang sama, dengan \code{NIM} ditulis secara konsisten (huruf besar/kecil mengikuti NIM resmi). Pastikan seluruh berkas berada tepat di dalam satu arsip \code{.zip} tanpa struktur folder bertingkat yang tidak perlu.

\begin{callout}[ctpBlue]
\textbf{\color{ctpBlue}Catatan Pengerjaan}
\begin{itemize}
  \item \textbf{Soal 1--4:} Wajib dikerjakan oleh semua mahasiswa. Soal 1 wajib dilaporkan lebih dulu sebagai baseline.
  \item \textbf{Soal 5:} Wajib dikerjakan; bagian (b) memerlukan kode, bagian (a) dan (c) bersifat analitis.
  \item \textbf{Encoder beku:} pengurangan akurasi akibat melatih encoder akan dinilai sebagai kesalahan prosedur.
  \item \textbf{Estimasi waktu:} 3--4 jam.
\end{itemize}
\end{callout}

\section*{Referensi}
{\small\color{ctpSubtext}
\begin{enumerate}[label={[\arabic*]}, leftmargin=2em]
  \item Baltrusaitis, T., Ahuja, C., \& Morency, L.-P. (2018). \emph{Multimodal Machine Learning: A Survey and Taxonomy}. IEEE TPAMI, 40(2), 421--443. --- taksonomi early/intermediate/late fusion.
  \item Arevalo, J., Solorio, T., Montes-y-Gomez, M., \& Gonzalez, F. A. (2017). \emph{Gated Multimodal Units for Information Fusion}. ICLR Workshop. --- dasar mekanisme \emph{gating}.
  \item Liang, P., Zadeh, A., \& Morency, L.-P. (2023). \emph{Tutorial on Multimodal Machine Learning}. ICML. --- bagian strategi fusi.
  \item Goodfellow, I., Bengio, Y., \& Courville, A. (2016). \emph{Deep Learning}, Bab 6--7: Deep Feedforward Networks. MIT Press.
\end{enumerate}}

\end{document}
